\section{Whole centres on covers disjoint only on the original cosupport}
\label{sec:hironaka-support-disjoint-canonical-centre}

\begin{theorem}[Exact whole centre traces on a support-disjoint cover]
\label{hir:support-disjoint-whole-canonical-centre}
\uses{hir:componentwise-canonical-centre-exact-trace,hir:open-restriction-marked,hir:invariant-driver,hir:centre-function-seq}
\lean{Hironaka.HironakaSupportDisjointCanonicalCentre.parallelCentre_subset_original_cosupp,Hironaka.HironakaSupportDisjointCanonicalCentre.preimage_parallelCentre_support_disjoint,Hironaka.HironakaSupportDisjointCanonicalCentre.parallelCentre_admissible_support_disjoint}
\leanok
Let $M$ be an original marked ideal in the domain of a proved
invariant $R$, put $S=\operatorname{cosupp}M$, and let $T:Y\to X$
be an actual test prefix. Let original opens $U_i$ cover $S$,
with the sets $U_i\cap S$ pairwise disjoint. Suppose each true
restricted prefix $T|_{U_i}$ is canonical for $M|_{U_i}$.
Form the union $C(T)$ of the whole local canonical maximum
centres in the actual stage preimages of these opens.
Then $C(T)\subset T^{-1}S$, its entire trace on each
$T^{-1}U_i$ is exactly that whole local maximum centre,
and $C(T)$ is a globally closed smooth admissible centre.
The opens themselves may overlap away from $S$.
\end{theorem}
\begin{proof}
A point of a local maximum belongs to its genuine local
current cosupport. The whole current marked-record restriction
dictionary identifies it with a point of the global current
cosupport. Its original image lies in $S$, since every
current cosupport maps into the original cosupport along
an actual test prefix. This proves $C(T)\subset T^{-1}S$.

If a point in $T^{-1}U_i$ lies in the union through the
centre on $T^{-1}U_j$, its original image lies in both
$U_i\cap S$ and $U_j\cap S$. Support disjointness forces
$i=j$; injectivity of the true stage-open inclusion then
identifies the point with the same whole local centre point.
The reverse trace inclusion is immediate.

The original cosupport $S$ is closed by coherence and
regularity of the original carrier. The opens $U_i$ together
with $X\setminus S$ cover $X$. On each actual stage preimage
of $U_i$ the exact trace is genuinely admissible by the
invariant's maximum-centre theorem and the whole-record
restriction identity. On the preimage of $X\setminus S$
the trace is empty and hence admissible. Admissibility
descends from this genuine open cover: local complements
give closedness, local flattenings give smoothness and
the same local coordinates preserve every ordered boundary
member. This proves global admissibility without cutting
off any centre.
\end{proof}

\begin{theorem}[Actual proper support-disjoint refinement]
\label{hir:support-disjoint-canonical-global-refinement}
\uses{hir:support-disjoint-whole-canonical-centre,hir:resolution-driver,hir:open-restriction-tower}
\lean{Hironaka.HironakaSupportDisjointCanonicalCentre.exists_support_disjoint_canonical_refinement}
\leanok
Under the preceding original and actual-prefix hypotheses,
there is a genuine proper global blowup along $C(T)$.
The enlarged global original prefix is a test prefix, and
its actual restriction to every $U_i$ is canonical for
$M|_{U_i}$, with every empty stage position retained.
\end{theorem}
\begin{proof}
The constructed global centre is admissible and therefore
lies in the current cosupport. The invariant-domain condition
persists along the actual test prefix, so that cosupport
has empty interior. Hence the centre also has empty interior.
The genuine blowup-existence theorem applies to this global
closed smooth centre; the resulting actual blowup is proper.
Admissibility makes the enlarged original tower a test prefix.
On each $U_i$ its true next centre is the entire inverse image
of $C(T)$, which the preceding exact trace identifies with
the whole local canonical maximum. Thus every true restricted
prefix is extended by its genuine whole-maximum step and
remains canonical, including when that maximum is empty.
\end{proof}

\begin{remark}
This generalizes disjointness of ambient opens to disjointness
only on the fixed ORIGINAL cosupport. It does not assume
compatibility of chosen stage centres: their exact traces are
proved using the actual projection to that same original
cosupport. No general coherent cosupport is asserted to admit
such a support-disjoint cover by compact-cosupport opens.
\end{remark}
